\documentclass[10pt,a4paper]{article}

\usepackage[utf8]{inputenc}
\usepackage[T1]{fontenc}
\usepackage{graphicx}
\usepackage{booktabs}
\usepackage{amsmath}
\usepackage{hyperref}
\usepackage{float}
\usepackage{geometry}
\usepackage{caption}
\usepackage{subcaption}
\usepackage{xcolor}
\usepackage{array}

\geometry{margin=1in}
\graphicspath{{./}}
\hypersetup{
    colorlinks=true,
    linkcolor=blue,
    citecolor=blue,
    urlcolor=blue
}

\title{\textbf{Explainable GNN for Research Group Prediction in the AIFB Dataset}}
\author{
Anmol Raj Srivastav [4064093] \and
Janhavi Palav [4063279] \\
University of Paderborn, Germany
}
\date{}

\begin{document}
\maketitle

\begin{abstract}
This mini-project studies explainable graph neural networks for node
classification on the AIFB RDF knowledge graph. The task is to predict the
research group affiliation of a person using graph structure, typed RDF
relations, and a relational graph neural network. We train a two-layer
FastRGCN model with learnable node embeddings and evaluate it on the standard
AIFB train/test split. The trained model reaches 100.00\% training accuracy and
94.44\% test accuracy, with a test Macro-F1 score of 0.9508. To make the model
more transparent, we apply GNNExplainer from PyTorch Geometric to produce both
global relation-level explanations and local person-level explanation
subgraphs. The explanations are evaluated quantitatively using Fidelity+,
Fidelity-, Sparsity, and Stability. The results show that the model achieves
strong predictive performance while the explanation masks are sparse and
stable, although some RDF literal labels remain noisy and require careful
interpretation.
\end{abstract}

\noindent\textbf{Keywords:} Explainable Artificial Intelligence, Graph Neural
Networks, RDF Knowledge Graphs, FastRGCN, GNNExplainer, AIFB

\section{Introduction}
Graph-structured data occurs naturally in knowledge graphs, social networks,
citation networks, and biological networks. In such data, entities are
represented as nodes and relationships are represented as edges. Standard
machine learning methods often treat observations as independent examples and
therefore do not directly exploit these relational dependencies. Graph Neural
Networks (GNNs) address this limitation by learning node representations through
message passing: a node representation is updated by aggregating information
from its neighbours.

This project focuses on an RDF knowledge graph, where relationships are not all
of the same type. A relation such as \texttt{author} has a different meaning
from a relation such as \texttt{publication}, \texttt{phone}, or
\texttt{worksAtProject}. For this reason, a relational GNN is more appropriate
than a plain homogeneous GCN. We use FastRGCN, an efficient implementation of
Relational Graph Convolutional Networks, to model the typed relations in the
AIFB graph.

The prediction task is to classify a person into one of four research groups:
Business Information and Communication Systems, Efficient Algorithms, Knowledge
Management, or Complexity Management. Since GNNs can be difficult to interpret,
we combine the FastRGCN classifier with GNNExplainer. The goal is not only to
obtain accurate predictions, but also to understand which edges and relation
types support those predictions.

\section{Data Analysis}
The AIFB dataset is an RDF dataset describing people, publications, research
groups, and other academic entities. RDF data is represented as triples of the
form subject--predicate--object. In graph terms, subjects and objects become
nodes, while predicates become typed edges. The AIFB task is a node
classification task on labelled person nodes.

The original graph contains direct target-label relations such as affiliation.
The processed PyTorch Geometric AIFB dataset removes target leakage and adds
inverse relations, allowing messages to flow in both directions. This produces
a graph with 8,285 nodes and 90 relation types.

\begin{table}[H]
\centering
\caption{AIFB dataset statistics used in this project.}
\label{tab:data-stats}
\begin{tabular}{lr}
\toprule
\textbf{Statistic} & \textbf{Value} \\
\midrule
Total RDF triples in full RDF file & 29,226 \\
Unique graph nodes after PyG processing & 8,285 \\
Relation types after adding inverse relations & 90 \\
Training person nodes & 140 \\
Test person nodes & 36 \\
Target classes & 4 \\
\bottomrule
\end{tabular}
\end{table}

\begin{figure}[H]
\centering
\includegraphics[width=0.78\textwidth]{class_distribution.png}
\caption{Distribution of persons across research group classes. The dataset is
imbalanced: Business Information and Communication Systems and Knowledge
Management contain more labelled persons than the other two classes.}
\label{fig:class-dist}
\end{figure}

\begin{figure}[H]
\centering
\includegraphics[width=0.82\textwidth]{top_predicates.png}
\caption{Most frequent predicates in the full AIFB RDF graph. Frequent
relations include publication and author information, which are useful
structural signals for predicting research group membership.}
\label{fig:top-predicates}
\end{figure}

The class distribution in Figure~\ref{fig:class-dist} shows that the dataset is
small and imbalanced. This is important when interpreting model performance:
accuracy alone can hide weaker performance on minority classes. For that
reason, we report both accuracy and Macro-F1.

\section{Model Training and Evaluation}
\subsection{Architecture}
The model used in this project is a two-layer FastRGCN. Since the AIFB dataset
does not provide dense node features, the model starts with a learnable node
embedding table. Each node id is mapped to a 32-dimensional embedding. The
embeddings are then processed by two FastRGCN convolution layers. The first
layer maps the embedding to a hidden representation with 16 channels, followed
by ReLU and dropout. The second layer maps hidden representations to four class
logits.

\begin{table}[H]
\centering
\caption{FastRGCN model configuration.}
\label{tab:model-config}
\begin{tabular}{ll}
\toprule
\textbf{Component} & \textbf{Setting} \\
\midrule
Input representation & Learnable node embedding \\
Embedding dimension & 32 \\
Hidden dimension & 16 \\
GNN layers & 2 FastRGCNConv layers \\
Output classes & 4 research groups \\
Dropout & 0.2 \\
Optimizer & Adam \\
Learning rate & 0.01 \\
Weight decay & 0.0005 \\
Epochs & 50 \\
Loss & Cross-entropy loss \\
\bottomrule
\end{tabular}
\end{table}

\subsection{Predictive Performance}
The model was trained for 50 epochs and evaluated on both the training and test
sets. The final performance is shown in Table~\ref{tab:model-performance}.

\begin{table}[H]
\centering
\caption{FastRGCN predictive performance on AIFB.}
\label{tab:model-performance}
\begin{tabular}{lcc}
\toprule
\textbf{Split} & \textbf{Accuracy} & \textbf{Macro-F1} \\
\midrule
Train & 100.00\% & 1.0000 \\
Test & 94.44\% & 0.9508 \\
\bottomrule
\end{tabular}
\end{table}

\begin{table}[H]
\centering
\caption{Per-class classification report on the test set.}
\label{tab:classification-report}
\begin{tabular}{lrrrr}
\toprule
\textbf{Class} & \textbf{Precision} & \textbf{Recall} &
\textbf{F1-score} & \textbf{Support} \\
\midrule
Business Information and Communication Systems & 0.88 & 1.00 & 0.94 & 15 \\
Efficient Algorithms & 1.00 & 0.92 & 0.96 & 12 \\
Knowledge Management & 1.00 & 1.00 & 1.00 & 3 \\
Complexity Management & 1.00 & 0.83 & 0.91 & 6 \\
\midrule
Macro average & 0.97 & 0.94 & 0.95 & 36 \\
Weighted average & 0.95 & 0.94 & 0.94 & 36 \\
\bottomrule
\end{tabular}
\end{table}

\begin{figure}[H]
\centering
\includegraphics[width=0.92\textwidth]{training_curves.png}
\caption{Training loss and accuracy curves over 50 epochs. The training
accuracy reaches 100\%, while the test accuracy remains high at 94.44\%.}
\label{fig:training-curves}
\end{figure}

The model performs well despite the small labelled dataset. The high training
accuracy indicates that the model can fit the labelled training nodes, while the
test accuracy and Macro-F1 show that the learned relational patterns generalise
reasonably well. The lower recall for Complexity Management is expected because
this class has fewer examples than the larger groups.

\section{Model Explanation}
\subsection{GNNExplainer Setup}
To explain the trained FastRGCN, we use GNNExplainer from PyTorch Geometric.
The explainer learns a soft mask over graph edges. Edges with high mask values
are interpreted as more important for preserving the model prediction.

The saved FastRGCN model uses a learnable embedding layer. For compatibility
with GNNExplainer, the embeddings are first computed once and then passed to a
wrapper model containing the same trained FastRGCN convolution layers. This
means that the explanation is still generated for the trained FastRGCN model,
not for a separate surrogate model.

\subsection{Global Explanation}
For the global explanation, GNNExplainer is applied to a sample of training
nodes. The edge-mask scores are averaged by relation type. The top relation
types indicate which relations are generally influential for model predictions.

\begin{figure}[H]
\centering
\includegraphics[width=0.86\textwidth]{global_explanation.png}
\caption{Global explanation showing the top relation types ranked by average
edge-mask importance across sampled training nodes.}
\label{fig:global-explanation}
\end{figure}

The global explanation suggests that metadata and structural relations such as
\texttt{phone}, \texttt{type}, \texttt{fax}, \texttt{head\_inverse}, and
project-related relations can influence the model. Some highly ranked relations
are not semantically ideal from a human perspective, which is an important
limitation of this dataset and preprocessing pipeline. Nevertheless, the result
is useful because it shows which relation types the trained model is sensitive
to.

\subsection{Local Person-Level Explanations}
Local explanations were generated for five correctly classified test persons.
For each person, the top-10 most important edges were extracted and visualised
with human-readable labels wherever possible. The person-level explanations are
also saved in a Markdown and CSV summary.

\begin{figure}[H]
\centering
\includegraphics[width=0.92\textwidth]{local_exp_Steffen_Lamparter.png}
\caption{Local explanation for Steffen Lamparter. The highlighted subgraph
shows the most influential edges used by the model to predict the class
Efficient Algorithms.}
\label{fig:local-steffen}
\end{figure}

\begin{table}[H]
\centering
\caption{Summary of local person explanations.}
\label{tab:local-summary}
\begin{tabular}{p{3.2cm}p{4.2cm}ccc}
\toprule
\textbf{Person} & \textbf{Predicted group} & \textbf{Confidence} &
\textbf{Correct} & \textbf{Important edges} \\
\midrule
Steffen Lamparter & Efficient Algorithms & 1.000 & Yes & 7 \\
Nenad Stojanovic & Efficient Algorithms & 1.000 & Yes & 17 \\
Sebastian Blohm & Efficient Algorithms & 1.000 & Yes & 2 \\
Yu Li & Business Information and Communication Systems & 1.000 & Yes & 6 \\
J\"urgen Angele & Efficient Algorithms & 1.000 & Yes & 7 \\
\bottomrule
\end{tabular}
\end{table}

For \textbf{Steffen Lamparter}, the model predicts the class Efficient
Algorithms with high confidence. The explanation shows that the model mainly
relies on project- and publication-related connections around the person node.
Important relations include \texttt{carriedOutBy},
\texttt{worksAtProject\_inverse}, and \texttt{member}. This suggests that the
model is not simply memorising the person node, but using nearby academic
activity such as projects, publications, and related collaborators to infer the
research group. The strongest evidence connects Steffen Lamparter to
project/publication entities that are also linked to other researchers,
indicating that shared research context is an important classification signal.

For \textbf{Nenad Stojanovic}, the model again predicts Efficient Algorithms.
The most influential edges include \texttt{member}, \texttt{carriedOutBy},
\texttt{financedBy\_inverse}, and \texttt{finances}. These relations connect
the person to organisational, project, and research activity nodes. The
explanation indicates that the model uses the surrounding collaboration and
project structure to support the prediction. In particular, the presence of
repeated links through project-like and member-like relations suggests that the
model identifies Nenad Stojanovic's research group from the local neighbourhood
pattern rather than from a single isolated edge.

For \textbf{Sebastian Blohm}, the model predicts Efficient Algorithms. The
explanation is more compact than for some other nodes, with only a small number
of highly important edges above the threshold. The strongest relation is
\texttt{worksAtProject\_inverse}, connecting Sebastian Blohm to a nearby
project or publication-related entity. This means the model's decision is
mainly driven by a small but strong local structure. The explanation suggests
that even a few project or publication links can be enough for the FastRGCN
model to identify the research group when those links are strongly associated
with a class-specific neighbourhood.

For \textbf{Yu Li}, the model predicts Business Information and Communication
Systems. The important edges include \texttt{member},
\texttt{publication\_inverse}, \texttt{publishes}, and \texttt{author}. This
explanation is intuitive because publication and authorship relations are
meaningful indicators of academic group membership. The model appears to use Yu
Li's connections to publications and author-related entities as evidence for
the predicted research group. Compared with metadata-only explanations, this
case is easier to interpret because the influential edges correspond directly
to academic activity.

For \textbf{J\"urgen Angele}, the model predicts Efficient Algorithms. The
explanation highlights relations such as \texttt{worksAtProject\_inverse},
\texttt{carriedOutBy}, \texttt{carriesOut\_inverse},
\texttt{isWorkedOnBy}, \texttt{publishes}, and \texttt{editor}. These
relations connect the person to projects, publications, and research outputs.
The model therefore seems to base its prediction on J\"urgen Angele's research
environment: the projects and publications around the person node form a local
pattern that is characteristic of the predicted group. This explanation
supports the idea that the GNN learns from relational neighbourhood structure
rather than only from direct labels.

Overall, the local explanations show that the FastRGCN model mainly relies on
project, publication, authorship, membership, and collaboration-related edges.
This is reasonable for the AIFB dataset because research group affiliation is
naturally reflected in a person's academic neighbourhood. Some RDF labels are
still noisy because RDF objects can be long abstracts, numbers, file-like
strings, or metadata values. However, the important relation types still show
that the model uses meaningful graph structure to make its predictions.

\section{Explanation Evaluation}
The explanations are evaluated using four metrics. Fidelity+ measures whether
the model prediction changes when important edges are removed. Fidelity-
measures whether the prediction stays the same when only important edges are
kept. Sparsity measures how compact the explanation is, and Stability measures
how similar repeated explanation masks are.

\begin{table}[H]
\centering
\caption{Quantitative evaluation of GNNExplainer over 20 test nodes.}
\label{tab:explanation-metrics}
\begin{tabular}{lcc}
\toprule
\textbf{Metric} & \textbf{Mean} & \textbf{Standard deviation} \\
\midrule
Fidelity+ & 0.2000 & 0.4104 \\
Fidelity- & 1.0000 & 0.0000 \\
Sparsity & 0.9996 & 0.0004 \\
Stability & 0.9989 & 0.0032 \\
\bottomrule
\end{tabular}
\end{table}

\begin{figure}[H]
\centering
\includegraphics[width=0.76\textwidth]{eval_metrics.png}
\caption{Explanation quality metrics averaged over 20 test nodes.}
\label{fig:eval-metrics}
\end{figure}

The high Fidelity- score indicates that the selected important edges are
sufficient to preserve the model prediction. The high Sparsity score shows that
only a very small fraction of the full graph is selected as important, so the
explanations are compact. Stability is also high, meaning repeated explainer
runs produce very similar edge masks. Fidelity+ is lower, which means that
removing the selected important edges does not always change the prediction.
This can happen when a graph contains redundant evidence paths: even after the
top explanation edges are removed, the model may still find alternative
relations that support the same class.

\section{Discussion}
The results show that FastRGCN is a suitable model for the AIFB RDF node
classification task. It explicitly handles typed relations and achieves strong
test performance. The explanation pipeline also satisfies the goal of connecting
model predictions to graph evidence. In particular, local explanations can be
shown as subgraphs with labelled nodes and labelled relation types.

At the same time, the explanations reveal several limitations. First, the AIFB
dataset is small, so the model can achieve high training accuracy quickly.
Second, some explanation edges correspond to metadata such as phone or fax
relations, which are not always semantically satisfying. Third, node labels
extracted from RDF are sometimes noisy. These limitations do not invalidate the
approach, but they should be considered when interpreting the explanation
figures.

\section{Conclusion}
This mini-project implemented an explainable GNN pipeline for the AIFB RDF
dataset. A FastRGCN model was trained to predict research group affiliations
for person nodes and achieved 94.44\% test accuracy with a Macro-F1 score of
0.9508. GNNExplainer was then used to generate global relation-level
explanations and local person-level subgraph explanations. Quantitative
evaluation showed that the explanations are sparse, stable, and sufficient to
preserve predictions when only important edges are retained. Overall, the
project demonstrates how relational GNNs and explanation methods can be
combined to make predictions on RDF knowledge graphs more transparent.

\section{Contributions}
Anmol Raj Srivastav worked on the implementation of the data analysis,
FastRGCN training pipeline, GNNExplainer integration, explanation
visualisation, and quantitative evaluation scripts. Janhavi Palav contributed to
report writing, result interpretation, and the organisation of the final
presentation material. Both team members reviewed the final results and helped
connect the implementation to the XAI mini-project requirements.

\begin{thebibliography}{9}

\bibitem{gnnexplainer}
Ying, R., Bourgeois, D., You, J., Zitnik, M., Leskovec, J.:
GNNExplainer: Generating Explanations for Graph Neural Networks.
In: Advances in Neural Information Processing Systems 32 (2019).

\bibitem{pyg}
Fey, M., Lenssen, J. E.:
Fast Graph Representation Learning with PyTorch Geometric.
In: ICLR Workshop on Representation Learning on Graphs and Manifolds (2019).

\bibitem{rgcn}
Schlichtkrull, M., Kipf, T. N., Bloem, P., van den Berg, R.,
Titov, I., Welling, M.:
Modeling Relational Data with Graph Convolutional Networks.
In: European Semantic Web Conference, pp. 593--607. Springer (2018).

\bibitem{aifb}
AIFB RDF Dataset.
\url{https://figshare.com/articles/dataset/AIFB_DataSet/745364/1}

\bibitem{pygentities}
PyTorch Geometric Entities Dataset Documentation.
\url{https://pytorch-geometric.readthedocs.io/}

\end{thebibliography}

\end{document}
